%===========================================================================


\section{Action-Derived Features: An Ablation \texorpdfstring{[E]}{[E]}}
\label{sec:ablation}
%===========================================================================

\pcAblNDropped{} of the \pcHTwoNTerms{} covariates record clinician actions:
vasopressor exposure and the per-test indicators of whether a laboratory value
was obtained in a given hour. Cross-label coefficient instability is
concentrated in these features (Section~\ref{sec:h2_action}). This section
tests whether the models depend on them by removing them and refitting, with
the specification otherwise unchanged (Section~\ref{sec:ablation_method}).

\begin{table}[H]
    \centering\small
    \adjustbox{max width=\textwidth}{%
    \begin{tabular}{llrrrrr}
        \toprule
        \textbf{Var} & \textbf{Model} & \textbf{AUROC full} & \textbf{AUROC ablated}
        & \textbf{$\Delta$} & \textbf{Utility$_{\max}$ full}
        & \textbf{Utility$_{\max}$ ablated} \\
        \midrule
        A & Primary & \pcAurocPrimaryA & \pcAblAurocPrimaryA & \pcAblDeltaAurocPrimaryA & \pcUtilityMaxPrimaryA & \pcAblUtilityMaxPrimaryA \\
        A & GBT     & \pcAurocGbtA     & \pcAblAurocGbtA     & \pcAblDeltaAurocGbtA     & \pcUtilityMaxGbtA     & \pcAblUtilityMaxGbtA \\
        \midrule
        B & Primary & \pcAurocPrimaryB & \pcAblAurocPrimaryB & \pcAblDeltaAurocPrimaryB & \pcUtilityMaxPrimaryB & \pcAblUtilityMaxPrimaryB \\
        B & GBT     & \pcAurocGbtB     & \pcAblAurocGbtB     & \pcAblDeltaAurocGbtB     & \pcUtilityMaxGbtB     & \pcAblUtilityMaxGbtB \\
        \midrule
        C & Primary & \pcAurocPrimaryC & \pcAblAurocPrimaryC & \pcAblDeltaAurocPrimaryC & \pcUtilityMaxPrimaryC & \pcAblUtilityMaxPrimaryC \\
        C & GBT     & \pcAurocGbtC     & \pcAblAurocGbtC     & \pcAblDeltaAurocGbtC     & \pcUtilityMaxGbtC     & \pcAblUtilityMaxGbtC \\
        \bottomrule
    \end{tabular}}
    \caption{Discrimination and best-threshold Utility Score with and without
    the \pcAblNDropped{} action-derived features. $\Delta$ is a paired
    difference; Table~\ref{tab:ablation_ci} carries its interval.
    \textbf{[E] exploratory.}}
    \label{tab:ablation}
\end{table}

\subsection{Coefficient Stability After Ablation}

Table~\ref{tab:ablation_stability} compares cross-label coefficient stability
before and after ablation.

\begin{table}[H]
    \centering
    \begin{tabular}{lrrr}
        \toprule
        \textbf{Specification} & \textbf{Covariates} & \textbf{Unstable}
        & \textbf{\% unstable} \\
        \midrule
        Full     & \pcHTwoNTerms     & \pcHTwoEst            & \pcHTwoPctUnstable \\
        Ablated  & \pcAblHTwoNTerms  & \pcAblHTwoNUnstable   & \pcAblHTwoPctUnstable \\
        \bottomrule
    \end{tabular}
    \caption{Cross-label coefficient instability before and after removing the
    action-derived features, under the identical stability rule
    (sign reversal or magnitude change $\geq 2\times$ across Variants A, B
    and C). \textbf{[E] exploratory.}}
    \label{tab:ablation_stability}
\end{table}

Removing \pcAblNDropped{} of the \pcHTwoNTerms{} covariates changes AUROC by
\pcAblDeltaAurocPrimaryB{} for the primary model (95\% CI
\pcAblDeltaCiPrimaryB) and \pcAblDeltaAurocGbtB{} for GBT (\pcAblDeltaCiGbtB)
under the primary label, and by no more than \pcAblDeltaAurocPrimaryC{} for
either model under any variant (Table~\ref{tab:ablation_ci}).
\pcAblNExcludeZero{} of the six paired contrasts exclude zero, so the cost is
small, below 0.01 AUROC everywhere, but not absent; no practical-equivalence
margin was pre-registered.

\begin{table}[H]
    \centering\small
    \begin{tabular}{llrrl}
        \toprule
        \textbf{Var} & \textbf{Model} & \textbf{AUROC full}
        & \textbf{$\Delta$ (ablated $-$ full)} & \textbf{95\% CI} \\
        \midrule
        A & Primary & \pcAurocPrimaryA & \pcAblDeltaAurocPrimaryA & \pcAblDeltaCiPrimaryA \\
        A & GBT     & \pcAurocGbtA     & \pcAblDeltaAurocGbtA     & \pcAblDeltaCiGbtA \\
        \midrule
        B & Primary & \pcAurocPrimaryB & \pcAblDeltaAurocPrimaryB & \pcAblDeltaCiPrimaryB \\
        B & GBT     & \pcAurocGbtB     & \pcAblDeltaAurocGbtB     & \pcAblDeltaCiGbtB \\
        \midrule
        C & Primary & \pcAurocPrimaryC & \pcAblDeltaAurocPrimaryC & \pcAblDeltaCiPrimaryC \\
        C & GBT     & \pcAurocGbtC     & \pcAblDeltaAurocGbtC     & \pcAblDeltaCiGbtC \\
        \bottomrule
    \end{tabular}
    \caption[Paired ablation contrasts with intervals]{Paired
    ablated-minus-full AUROC contrasts with 95\% stay-level cluster-bootstrap
    intervals (Protocol Amendment 10), differenced within replicate.
    \textbf{[E] exploratory; enters no corrected family.}}
    \label{tab:ablation_ci}
\end{table}

The cross-label unstable count falls from \pcHTwoEst{} of \pcHTwoNTerms{}
covariates to \pcAblHTwoNUnstable{} of \pcAblHTwoNTerms{}
(\pcAblHTwoPctUnstable\%). This is not a subset relation: the single
remaining flag falls on \texttt{lactate}, which the full fit does not flag,
because the refit re-estimates the remaining coefficients. Together with the
by-kind split of Section~\ref{sec:h2_action} (\pcHTwoPctUnstableAction\%
unstable among action-derived features against \pcHTwoPctUnstablePhysio\% among
physiological ones), these results indicate that the features recording
clinician behaviour contribute little to discrimination and carry most of the
cross-label instability (Section~\ref{sec:ablation_establishes}).
